\section{Related work} \label{sec:related}

The logic $\GPtwo$ was shown decidable in \cite{percentage,twovarpres}, while the decidability
of $\Ctwo$ with cardinality constraints is implicit in \cite{tonyeryk, ultimatelyperiodic}. Our complexity bound for $\GPtwo$ refines the analysis of the Kleene star operator in \cite{usicalp}, which in turn relies on \cite{kleenestar}. Our analysis of $\Ctwo$ with constraints obviously relies heavily on the techniques developed by Pratt-Hartmann for analysis of $\Ctwo$, in particular normal forms, the categorization of types, and the distinction between large and small number of realizers (e.g. the notion of Z-differentiated \cite{ctwocomplexity,gctwocomplexity}).

Although we borrow techniques from prior work, our approach gives a self-contained proof of decidability for $\Ctwo$ itself, and also for $\Ctwo$ with global constraints in extensions of Presburger arithmetic. An alternative approach to get bounds for $\Ctwo$ with Presburger global constraints is by reducing to decidability of $\Ctwo$.
This approach was presented by Rudolph in \cite{rudolphcardinality}, relying on ideas from \cite{baadercardinality}. Although technically speaking the results in \cite{rudolphcardinality} are presented for description logics, they could easily be applied to give the same bounds for $\Ctwo$.

In addition to contributing tight bounds for complexity for decidability of $\GPtwo$ a well as the size of the spectrum  for $\Ctwo$, we note that \cite{usicalp} is phrased in terms of graph neural networks, while from the analysis of spectra in \cite{tonyeryk,ultimatelyperiodic} it would be difficult to see how semilinear representations are derived. Thus we believe our work additionally makes the analyses of logics with arithmetic substantially more accessible.
